\section*{Chapter~6~~Conclusion}
\setcounter{section}{6} \setcounter{subsection}{0} 
\setcounter{table}{0} \setcounter{figure}{0} \setcounter{equation}{0}\setcounter{definition}{0}
\addcontentsline{toc}{section}{\textbf{Chapter~6~~Conclusion}}

\vspace{4pt}

\noindent In this project, we developed the \textit{\textbf{BLEGuard}} system, a novel hybrid detection mechanism designed to safeguard Bluetooth Low Energy (BLE) networks against sophisticated spoofing attacks. BLEGuard's unique integration of a pre-detection scheme, reconstruction techniques, and classification models enables it to effectively identify and neutralize threats, thereby enhancing network security. The system's high detection accuracy, combined with a low false alarm rate and un-detection rate, underscores its potential not only as a specialized tool for BLE security but also for broader applications in industry, healthcare, and smart home sectors. The practical application of BLEGuard in these sectors can significantly mitigate risks associated with the inherent security vulnerabilities of BLE technologies, providing a reliable security solution that aligns with the needs of modern connected environments.

\vspace{12pt}

The construction of the \textit{\textbf{BLE-SAD}} dataset was a cornerstone of this project, providing a critical resource for testing and refining the BLEGuard system. This large-scale dataset, generated through comprehensive simulations of spoofing attacks within a controlled testbed environment, offers an invaluable asset for related cybersecurity research. Its detailed representation of varied attack scenarios enables researchers and security professionals to rigorously evaluate Bluetooth security solutions and contributes to the ongoing development of advanced defensive methods against an array of cyber threats.

\vspace{12pt}

Looking to the future, there are several avenues for further enhancing the BLEGuard system and the BLE-SAD dataset. Enriching the dataset with a broader spectrum of attack scenarios, including Man in the Middle (MITM) and Distributed Denial of Service (DDoS) attacks, will extend our research's reach and enrich the training resources available for developing robust BLE security measures. Additionally, adapting BLEGuard to function as a general framework for Bluetooth security could transform it into a versatile tool capable of defending against a diverse range of cyber threats. Such advancements will not only fortify the security of BLE networks but also pave the way for next-generation protection mechanisms in the evolving landscape of digital communication technologies.

\newpage
